\clearpage\begingroup\NoSourceHead\SourceLeaf{554}
\begin{center}
\setlength{\fboxsep}{1em}
\fbox{\begin{minipage}[c][7.6in][s]{.88\textwidth}
\hrule\vspace{1.8em}\centering
{\Large CARL FRIEDRICH GAUSS\par}
\vspace{2em}{\Huge WERKE\par}
\vspace{2em}ZEHNTEN BANDES ZWEITE ABTEILUNG\par
\vspace{1em}ABHANDLUNG 6\par
\vspace{4em}HERAUSGEGEBEN\par
\vspace{.8em}VON DER\par
\vspace{.8em}GESELLSCHAFT DER WISSENSCHAFTEN\par
\vspace{.8em}ZU\par
\vspace{.8em}{\large GÖTTINGEN\par}
\vfill{\small IN KOMMISSION BEI JULIUS SPRINGER IN BERLIN\par}
\vspace{.8em}1930\par
\end{minipage}}
\end{center}
\clearpage\NoSourceHead\SourceLeaf{556}
\begin{center}
\vspace*{.7in}{\Large GAUSS ALS ZAHLENRECHNER}\\[.8in]
VON\\[.8in]
{\large PHILIPP MAENNCHEN}
\end{center}
\vfill
\noindent\rule{\textwidth}{.3pt}\par
\begin{center}\small
Umgearbeiteter Abdruck des Aufsatzes: \textit{Die Wechselwirkung zwischen Zahlenrechnen und Zahlentheorie bei C. F. Gauss, Materialien für eine wissenschaftliche Biographie von Gauss}, gesammelt von \textsc{F. Klein, M. Brendel} und \textsc{L. Schlesinger}, Heft VII, 1918; Nachrichten von der Königl. Gesellschaft der Wissenschaften zu Göttingen, Mathematisch-physikalische Klasse, 1918.\\
Vorgelegt in der Sitzung vom 8. Februar 1918.
\end{center}
\noindent\rule{\textwidth}{.3pt}\par
\vspace{.7in}\noindent{\small X\textsubscript{2} Abh. 6.}\hfill{\small 1}\par
\clearpage\NoSourceHead\SourceLeaf{558}
\renewcommand{\thefootnote}{\arabic{footnote})}\setcounter{footnote}{0}
\begin{center}{\large\bfseries Einleitendes.}\end{center}
\noindent\hspace*{\parindent}Wenn man die Lebensarbeit des princeps mathematicorum überblickt, so
fällt neben der ausserordentlichen Vielseitigkeit und Tiefe wohl am meisten
die Tatsache auf, dass sich sehr viele Ergebnisse vorfinden, die durch zahlenmässiges Rechnen — und dazu häufig durch peinlich genaues Rechnen —
gewonnen worden sind. Und zwar sind es nicht nur Beispiele, die gewissermassen als Erläuterungen zu den von \textsc{Gauss} entdeckten allgemeinen Sätzen
dienen, sondern wir finden ganze Tafeln, deren Herstellung allein die Lebensarbeit manches Rechners vom gewöhnlichen Schlage ausfüllen würde. Dazu
kommt, dass diese Rechnungen und Tafeln auf Stellenzahlen ausgeführt sind,
die das in der Praxis erforderliche Mass weit übersteigen. Wir erwähnen
vorerst nur die Herstellung von Logarithmentafeln, die vollständige Dezimalbruchentwicklung der reziproken Werte aller Primzahlen und Primzahlpotenzen
innerhalb des ersten Tausenders\footnote{Werke II, S. 411—434.}, die Tafel der arithmetisch-geometrischen
Mittel\footnote{Werke III, S. 403.}, die Tabellen zur Cyklotechnie\footnote{Werke II, S. 477—495.} und das Bruchstück: \textit{Quadratorum
Myrias prima}\footnote{Werke II, S. 504—505.}.

Wie kommt es nun, dass \textsc{Gauss} nicht nur alle diese anscheinend so mühseligen und geisttötenden Rechnungen auszuführen vermochte, sondern dass
er auch noch Zeit übrig behielt zu seinen tiefdringenden Untersuchungen, zu
seiner praktischen Betätigung als Astronom, Geodät und Physiker, ja, dass
er sich sogar noch mit dem Versicherungswesen beschäftigen konnte, alles
Gegenstände, die wiederum einen ungeheuern Aufwand von Zahlenrechnung
nötig machten? Die Erklärung kann nur die folgende sein: Die Rechnungen
\par\smallskip\noindent\hfill{\small 1\textsuperscript{*}}\par
\clearpage\SourceLeaf{559}\ArithmeticHead{4}{L}{PHILIPP MAENNCHEN, GAUSS ALS ZAHLENRECHNER.}
\noindent waren für ihn weder mühselig, noch geisttötend, sondern \textsc{Gauss} benutzte, wo
es irgend anging, Vorteile, »artificia«, die sich ihm namentlich bei der Aufstellung von Tafeln sehr bald darboten und ihm die Fortführung der betreffenden Tabelle bedeutend erleichterten. Diese Vorteile erwiesen sich sehr
bald selbst wieder als Sonderfälle neuer zahlentheoretischer Sätze, und so gewann er beständig neues Material zu seinen Forschungen in der niederen und
höheren Zahlentheorie und zugleich immer wieder neue Mittel, um Zahlenrechnungen einfach und eigenartig zu erledigen.

Der Zweck der vorliegenden Arbeit ist, den Gedankengängen nachzuspüren, die sich bei der Rechenarbeit des Zahlengewaltigen abspielten, die
Rechenvorteile aufzudecken und zu begründen, die ihm ermöglichten, das
spielend zu bewältigen, was anderen Rechnern nur Mühe und Arbeit, aber
gewiss keine Freude gemacht hätte. Da nun häufig die Spuren der ursprünglichen Gedankengänge umsomehr verschwinden, je druckreifer eine Arbeit
wird, so mussten die Belege an der Urquelle gesucht werden, im handschriftlichen Nachlass. Dabei bestand zugleich die Absicht, die Entwicklungsstufen
festzustellen, die \textsc{Gauss}’ Rechenfertigkeit nach und nach durchlaufen hat.
Allein dieser Plan konnte nicht streng durchgeführt werden und zwar aus
zwei Gründen. Einmal zeigt \textsc{Gauss} als echtes Genie vieles Sprunghafte in
seiner Entwicklung, und dann wird die Sichtung des handschriftlichen Nachlasses nach entwicklungsgeschichtlichen Gesichtspunkten dadurch erschwert,
dass viele Blätter zweifellos in verschiedenen Lebensabschnitten des Meisters
immer wieder von neuem beschrieben worden sind, so dass die zeitliche Einordnung in vielen Fällen nicht gelingt. Immerhin war es möglich, auch nach
dieser Richtung hin einige Klarheit zu schaffen.
\par\clearpage\endgroup
